\subsection{Hạn chế hiện tại}

Mặc dù hệ thống đã nối được ba mô-đun chính, phiên bản hiện tại vẫn còn một số hạn chế cần được nhìn nhận thẳng thắn.

\begin{itemize}
    \item Mô hình nhận diện đang dùng bộ ASL-250 pre-trained, chưa phải mô hình VSL thuần được huấn luyện trên dữ liệu người dùng Việt Nam.
    \item Chất lượng nhận diện còn nhạy với ánh sáng, góc quay, khoảng cách tới camera, che khuất tay và tốc độ ký hiệu.
    \item Pipeline hiện nhận diện theo token/gloss có thao tác xác nhận, chưa xử lý hoàn toàn tự động câu ký hiệu liên tục dài.
    \item Sai số ở tầng nhận diện có thể lan sang tầng khôi phục câu và tầng dịch, đặc biệt khi một token quan trọng bị nhận sai.
    \item Metric của mô-đun tiếng Việt rất cao trên bộ dữ liệu có cấu trúc/template, nhưng chưa đủ để khẳng định chất lượng tương đương trên dữ liệu giao tiếp tự do.
    \item Mô-đun dịch Việt -- Lào đã có Train loss, Validation loss và BLEU trong \texttt{vilo.pdf}, nhưng chưa có chrF và TER truy vết trong workspace nên cần tái lập thêm nếu muốn đánh giá đầy đủ hơn.
    \item Hệ thống hiện mới dừng ở mức nguyên mẫu demo, chưa được đóng gói thành ứng dụng desktop/web/mobile hoàn chỉnh.
\end{itemize}

